\subsection*{The Underlying Set}

\begin{tcolorbox}[colback=defbox, colframe=defborder, arc=2pt,
  left=6pt, right=6pt, top=4pt, bottom=4pt,
  title={\small\textbf{Definition (Whole Numbers)}},
  fonttitle=\small\bfseries]
\begin{definition}[Whole Numbers]
\label{def:whole-numbers}
Let $\mathbb{N}$ be the one-based natural-number system developed in the
preceding chapter. Let $0$ be a new element with $0\notin\mathbb{N}$. Define
\[
\mathbb{W}\coloneqq \mathbb{N}\cup\{0\}.
\]
\end{definition}
\end{tcolorbox}

\begin{remark*}[Interpretation]
Elements of $\mathbb{W}$ are either the new element $0$ or positive natural
numbers. The natural numbers embed into $\mathbb{W}$ by the inclusion map
$n\mapsto n$. The element $0$ is not a successor-stage element of the original
one-based Peano system.
\end{remark*}

\begin{dependencies}
\begin{itemize}
  \item \hyperref[def:peano-system]{Peano System}
\end{itemize}
\end{dependencies}
